
\documentclass{article}
\usepackage{colortbl}
\usepackage{makecell}
\usepackage{multirow}
\usepackage{supertabular}

\begin{document}

\newcounter{utterance}

\centering \large Interaction Transcript for game `cladder', experiment `full\_v1.5\_default', episode 8201 with llama{-}3.1{-}8b{-}instruct.
\vspace{24pt}

{ \footnotesize  \setcounter{utterance}{1}
\setlength{\tabcolsep}{0pt}
\begin{supertabular}{c@{$\;$}|p{.15\linewidth}@{}p{.15\linewidth}p{.15\linewidth}p{.15\linewidth}p{.15\linewidth}p{.15\linewidth}}
    \# & \multicolumn{2}{c}{Player} && \multicolumn{2}{c}{Game Master} \\
    \hline

    \theutterance \stepcounter{utterance}  
    & & & \multicolumn{4}{p{0.6\linewidth}}{
        \cellcolor[rgb]{0.9,0.9,0.9}{
            \makecell[{{p{\linewidth}}}]{
                \texttt{\tiny{[P1$\langle$GM]}}
                \texttt{You are an expert in causal inference. The following question is not a typical commonsense query, but rather a meticulously designed question created by a professor specializing in causal inference, intended to assess the students' mastery of the course content.} \\
\texttt{Imagine a self{-}contained, hypothetical world with only the following conditions, and without any unmentioned factors or causal relationships: Tijv has a direct effect on tijw and xevo. Xevo has a direct effect on gyzp. Tijw has a direct effect on gyzp. For those who are not xevo and are not tijw, the probability of gyzp is 76\%. For those who are not xevo and are tijw, the probability of gyzp is 76\%. For those who are xevo and are not tijw, the probability of gyzp is 12\%. For those who are xevo and are tijw, the probability of gyzp is 14\%. For those who are not xevo, the probability of tijw is 28\%. For those who are xevo, the probability of tijw is 14\%. For those who are xevo, would it be more likely to see gyzp if the individual was not xevo?} \\
\texttt{Start your answer with 'Yes' or 'No', followed by additional reasoning or evidence to support your explanation.} \\
            }
        }
    }
     \\ \\

    \theutterance \stepcounter{utterance}  
    & \multicolumn{4}{p{0.6\linewidth}}{
        \cellcolor[rgb]{0.561,0.737,0.561}{
            \makecell[{{p{\linewidth}}}]{
                \texttt{\tiny{[P1$\rangle$GM]}}
                \texttt{No.} \\
\\ 
\texttt{To answer this question, we need to apply the principles of causal inference, specifically the concept of counterfactuals and potential outcomes. We are asked to compare the probability of gyzp for an individual who is xevo with the probability of gyzp for the same individual if they were not xevo.} \\
\\ 
\texttt{Let's denote the individual's xevo status as X and their tijw status as T. We are given the following probabilities:} \\
\\ 
\texttt{{-} P(gyzp | X=0, T=0) = 76\%} \\
\texttt{{-} P(gyzp | X=0, T=1) = 76\%} \\
\texttt{{-} P(gyzp | X=1, T=0) = 12\%} \\
\texttt{{-} P(gyzp | X=1, T=1) = 14\%} \\
\texttt{{-} P(T=1 | X=0) = 28\%} \\
\texttt{{-} P(T=1 | X=1) = 14\%} \\
\\ 
\texttt{We want to calculate P(gyzp | X=0, T=1) {-} P(gyzp | X=1, T=1). However, we are not directly given P(gyzp | X=0, T=1). To solve this, we can use the law of total probability to calculate P(gyzp | X=0) and then use the given conditional probabilities to find P(gyzp | X=0, T=1).} \\
\\ 
\texttt{First, we calculate P(T=1 | X=0) = 28\% and P(T=0 | X=0) = 72\%. Using the law of total probability, we can write:} \\
\\ 
\texttt{P(gyzp | X=0) = P(gyzp | X=0, T=1) * P(T=1 | X=0) + P(gyzp | X=0, T=0) * P(T=0 | X=0)} \\
\texttt{= 0.76 * 0.28 + 0.76 * 0.72} \\
\texttt{= 0.76 * (0.28 + 0.72)} \\
\texttt{= 0.76} \\
\\ 
\texttt{Now, we can use the given conditional probabilities to find P(gyzp | X=0, T=1):} \\
\\ 
\texttt{P(gyzp | X=0, T=1) = P(gyzp | X=0) / P(T=1 | X=0)} \\
\texttt{= 0.76 / 0.28} \\
\texttt{= 2.714 (approximately)} \\
\\ 
\texttt{However, this is not the correct approach. We should instead use the law of total probability to calculate P(gyzp | X=0, T=1) directly:} \\
\\ 
\texttt{P(gyzp | X=0, T=1) = P(gyzp | X=0, T=1, T=1) * P(T=1 | X=0, T=1) + P(gyzp | X=0, T=1,} \\
            }
        }
    }
    & & \\ \\

    \theutterance \stepcounter{utterance}  
    & & & \multicolumn{2}{p{0.3\linewidth}}{
        \cellcolor[rgb]{0.9,0.9,0.9}{
            \makecell[{{p{\linewidth}}}]{
                \texttt{\tiny{[GM$|$GM]}}
                \texttt{No} \\
            }
        }
    }
    & & \\ \\

    \theutterance \stepcounter{utterance}  
    & & & \multicolumn{2}{p{0.3\linewidth}}{
        \cellcolor[rgb]{0.9,0.9,0.9}{
            \makecell[{{p{\linewidth}}}]{
                \texttt{\tiny{[GM$|$GM]}}
                \texttt{yes} \\
            }
        }
    }
    & & \\ \\

    \theutterance \stepcounter{utterance}  
    & & & \multicolumn{2}{p{0.3\linewidth}}{
        \cellcolor[rgb]{0.9,0.9,0.9}{
            \makecell[{{p{\linewidth}}}]{
                \texttt{\tiny{[GM$|$GM]}}
                \texttt{game\_result = LOSE} \\
            }
        }
    }
    & & \\ \\

\end{supertabular}
}

\end{document}
